% figs.tex —— 由 ch3/make_figs.py 產生，請勿手動修改。
% 重新產生：python ch3/make_figs.py（第 3 章，動畫 W301–W308）

\newcommand{\ThreeTabSim}{%
\begin{tabular}{@{}lrrrrc@{}}
\toprule
 & \multicolumn{2}{c}{100 步內回過家} & \multicolumn{2}{c}{1000 步內回過家} & 走無窮多步\\
\cmidrule(lr){2-3}\cmidrule(lr){4-5}
格子 & 模擬 & 理論 & 模擬 & 理論 & 回家的機率\\
\midrule
一維（直線） & 930 & 920.4 & 977 & 974.8 & 1\\
二維（正方格子） & 579 & 583.2 & 676 & 678.2 & 1\\
三維（簡單立方） & 300 & 311.9 & 325 & 331.5 & 0.3405\\
\bottomrule
\end{tabular}}

\newcommand{\ThreeFigVirial}{%
\begin{tikzpicture}[>=Stealth,walk/.style={decorate,decoration={snake,amplitude=0.7pt,segment length=5pt},cWalk,thick},dot/.style={circle,fill=black,inner sep=1.1pt}]
\begin{scope}[xshift=0.0cm]
\node[dot,label=below:{\scriptsize$\vec x$}] (x) at (0,0) {};
\node[dot,label=below:{\scriptsize$\vec y$}] (y) at (3,0) {};
\coordinate (w) at (1.5,0.9);
\draw[walk] (x) .. controls (0.4,0.9) .. (w);
\draw[walk] (w) .. controls (2.6,0.9) .. (y);
\node[circle,fill=cWrong,inner sep=2.6pt] at (w) {};
\node[font=\small] at (1.5,-0.75) {$-t$};
\end{scope}
\begin{scope}[xshift=4.2cm]
\node[dot,label=below:{\scriptsize$\vec x$}] (x) at (0,0) {};
\node[dot,label=below:{\scriptsize$\vec y$}] (y) at (3,0) {};
\coordinate (w) at (1.5,0.9);
\draw[walk] (x) .. controls (0.4,0.9) .. (w);
\draw[walk] (w) .. controls (2.6,0.9) .. (y);
\draw[walk] (w) .. controls +(60:1.1) and +(120:1.1) .. (w);
\node[circle,fill=cWrong,inner sep=2.6pt] at (w) {};
\node[font=\small] at (1.5,-0.75) {$+t^2$};
\end{scope}
\begin{scope}[xshift=8.4cm]
\node[dot,label=below:{\scriptsize$\vec x$}] (x) at (0,0) {};
\node[dot,label=below:{\scriptsize$\vec y$}] (y) at (3,0) {};
\coordinate (w) at (1.5,0.9);
\draw[walk] (x) .. controls (0.4,0.9) .. (w);
\draw[walk] (w) .. controls (2.6,0.9) .. (y);
\draw[walk] (w) .. controls +(30:1.1) and +(90:1.1) .. (w);
\draw[walk] (w) .. controls +(90:1.1) and +(150:1.1) .. (w);
\node[circle,fill=cWrong,inner sep=2.6pt] at (w) {};
\node[font=\small] at (1.5,-0.75) {$-t^3$};
\end{scope}
\node[font=\small] at (12.6,0.3) {$+\cdots$};
\end{tikzpicture}}

\newcommand{\ThreeTabNever}{%
{\small\setlength{\tabcolsep}{3.6pt}\begin{tabular}{@{}lrrrrrrrrrrr@{}}
\toprule
$N$ & 0 & 1 & 2 & 3 & 4 & 5 & 6 & 7 & 8 & 9 & 10\\
\midrule
從不回家 & 1 & 2 & 2 & 4 & 6 & 12 & 20 & 40 & 70 & 140 & 252\\
全部 $2^N$ & 1 & 2 & 4 & 8 & 16 & 32 & 64 & 128 & 256 & 512 & 1\,024\\
$P_N$ & 1.000 & 1.000 & 0.500 & 0.500 & 0.375 & 0.375 & 0.312 & 0.312 & 0.273 & 0.273 & 0.246\\
第 $N$ 步在家 & 1 & -- & 2 & -- & 6 & -- & 20 & -- & 70 & -- & 252\\
\bottomrule
\end{tabular}}}

\newcommand{\ThreeTabNoReturn}{%
{\small\setlength{\tabcolsep}{4pt}\begin{tabular}{@{}rcccccccc@{}}
\toprule
 & \multicolumn{2}{c}{一維} & \multicolumn{3}{c}{二維（正方）} & \multicolumn{2}{c}{三維（sc）}\\
\cmidrule(lr){2-3}\cmidrule(lr){4-6}\cmidrule(lr){7-8}
$N$ & 精確 & $\sqrt{2/\pi N}$ & 精確 & $\pi/\ln N$ & $\frac{\pi}{\ln N+\ln8+\gamma}$ & 精確 & $0.659463+\frac{0.28696}{\sqrt N}$\\
\midrule
$10^{1}$ & 0.2461 & 0.2523 & 0.5792 & 1.3644 & 0.6335 & 0.74723 & 0.75021\\
$10^{2}$ & 0.0796 & 0.0798 & 0.4168 & 0.6822 & 0.4326 & 0.68806 & 0.68816\\
$10^{3}$ & 0.0252 & 0.0252 & 0.3218 & 0.4548 & 0.3285 & 0.66853 & 0.66854\\
$10^{4}$ & 0.0080 & 0.0080 & 0.2613 & 0.3411 & 0.2647 & 0.66233 & 0.66233\\
\bottomrule
\end{tabular}}}

\newcommand{\ThreeFigTail}{%
\begin{tikzpicture}[x=6.0cm,y=0.54cm,>=Stealth]
\fill[cExact!15] (0,0) rectangle (1,1);
\draw[->] (0,0) -- (1.1,0) node[right,font=\scriptsize] {$q$};
\draw[->] (0,0) -- (0,5.6) node[above,font=\scriptsize] {$q^{d-1}/q^{2}$};
\draw[cBar,thick] plot coordinates {(0.4472,5.0000) (0.4564,4.8002) (0.4656,4.6121) (0.4749,4.4349) (0.4841,4.2677) (0.4933,4.1097) (0.5025,3.9604) (0.5117,3.8191) (0.5209,3.6852) (0.5301,3.5582) (0.5393,3.4377) (0.5486,3.3232) (0.5578,3.2143) (0.5670,3.1107) (0.5762,3.0120) (0.5854,2.9180) (0.5946,2.8282) (0.6038,2.7426) (0.6130,2.6608) (0.6223,2.5826) (0.6315,2.5078) (0.6407,2.4362) (0.6499,2.3676) (0.6591,2.3019) (0.6683,2.2388) (0.6775,2.1784) (0.6868,2.1203) (0.6960,2.0645) (0.7052,2.0109) (0.7144,1.9594) (0.7236,1.9098) (0.7328,1.8621) (0.7420,1.8162) (0.7512,1.7719) (0.7605,1.7292) (0.7697,1.6881) (0.7789,1.6484) (0.7881,1.6100) (0.7973,1.5731) (0.8065,1.5373) (0.8157,1.5028) (0.8250,1.4694) (0.8342,1.4371) (0.8434,1.4059) (0.8526,1.3757) (0.8618,1.3464) (0.8710,1.3181) (0.8802,1.2906) (0.8894,1.2640) (0.8987,1.2383) (0.9079,1.2133) (0.9171,1.1890) (0.9263,1.1655) (0.9355,1.1426) (0.9447,1.1205) (0.9539,1.0989) (0.9631,1.0780) (0.9724,1.0577) (0.9816,1.0379) (0.9908,1.0187) (1.0000,1.0000)};
\draw[cGauss,thick] plot coordinates {(0.2000,5.0000) (0.2133,4.6875) (0.2267,4.4118) (0.2400,4.1667) (0.2533,3.9474) (0.2667,3.7500) (0.2800,3.5714) (0.2933,3.4091) (0.3067,3.2609) (0.3200,3.1250) (0.3333,3.0000) (0.3467,2.8846) (0.3600,2.7778) (0.3733,2.6786) (0.3867,2.5862) (0.4000,2.5000) (0.4133,2.4194) (0.4267,2.3438) (0.4400,2.2727) (0.4533,2.2059) (0.4667,2.1429) (0.4800,2.0833) (0.4933,2.0270) (0.5067,1.9737) (0.5200,1.9231) (0.5333,1.8750) (0.5467,1.8293) (0.5600,1.7857) (0.5733,1.7442) (0.5867,1.7045) (0.6000,1.6667) (0.6133,1.6304) (0.6267,1.5957) (0.6400,1.5625) (0.6533,1.5306) (0.6667,1.5000) (0.6800,1.4706) (0.6933,1.4423) (0.7067,1.4151) (0.7200,1.3889) (0.7333,1.3636) (0.7467,1.3393) (0.7600,1.3158) (0.7733,1.2931) (0.7867,1.2712) (0.8000,1.2500) (0.8133,1.2295) (0.8267,1.2097) (0.8400,1.1905) (0.8533,1.1719) (0.8667,1.1538) (0.8800,1.1364) (0.8933,1.1194) (0.9067,1.1029) (0.9200,1.0870) (0.9333,1.0714) (0.9467,1.0563) (0.9600,1.0417) (0.9733,1.0274) (0.9867,1.0135) (1.0000,1.0000)};
\draw[cExact,thick] (0,1) -- (1,1);
\draw[cBar,->] (0.4472,5) -- (0.4472,5.5); \draw[cGauss,->] (0.2,5) -- (0.2,5.5);
\node[cBar,font=\scriptsize,right] at (0.56,5.0) {$d=1$：$1/q^2$，面積 $\to\infty$};
\node[cGauss,font=\scriptsize,right] at (0.56,4.3) {$d=2$：$1/q$，面積 $\to\infty$（對數）};
\node[cExact,font=\scriptsize,right] at (0.56,3.6) {$d=3$：$1$，面積有限（綠色）};
\foreach \x in {0.5,1} \node[below,font=\scriptsize] at (\x,0) {\x};
\node[below,font=\scriptsize] at (0,0) {0};
\foreach \v in {1,2,3,4,5} {\draw (-0.012,\v) -- (0.012,\v); \node[left,font=\scriptsize] at (-0.012,\v) {\v};}
\end{tikzpicture}}

\newcommand{\ThreeTabGzero}{%
\begin{tabular}{@{}lcccc@{}}
\toprule
$G(z;0)$，$z=x\,z_c$ & $x=0.9$ & $0.99$ & $0.999$ & $1$\\
\midrule
一維 $1/\sqrt{1-x^2}$ & 2.294 & 7.089 & 22.366 & $\infty$\\
正方 $\frac2\pi K(x)$ & 1.452 & 2.137 & 2.862 & $\infty$\\
sc $\int_0^\infty e^{-t}I_0(xt/3)^3dt$ & 1.233 & 1.409 & 1.480 & 1.516\\
\bottomrule
\end{tabular}}

\newcommand{\ThreeTabLattices}{%
{\small\begin{tabular}{@{}lclcccc@{}}
\toprule
格子 & 鄰居 & $G(z_c;0)$ 的閉式 & 數值 & 逃脫 $1/G$ & 回家 $1-1/G$ & 平均回家次數 $G-1$\\
\midrule
sc & 6 & $\frac{\sqrt6}{32\pi^3}\Gamma(\frac1{24})\Gamma(\frac5{24})\Gamma(\frac7{24})\Gamma(\frac{11}{24})$ & 1.516386 & 0.659463 & 0.340537 & 0.516386\\
bcc & 8 & $\frac{\Gamma(1/4)^4}{4\pi^3}$ & 1.393204 & 0.717770 & 0.282230 & 0.393204\\
fcc & 12 & $\frac{9\,\Gamma(1/3)^6}{2^{14/3}\pi^4}$ & 1.344661 & 0.743682 & 0.256318 & 0.344661\\
\bottomrule
\end{tabular}}}

\newcommand{\ThreeTabFm}{%
\begin{tabular}{@{}lccccc@{}}
\toprule
回家 $m$ 次 & $0$ & $1$ & $2$ & $3$ & $\ge4$\\
\midrule
$f_m=R^m(1-R)$（走無窮多步） & 0.6595 & 0.2246 & 0.0765 & 0.0260 & 0.0134\\
1000 步內的精確值 & 0.6685 & 0.2217 & 0.0735 & 0.0243 & 0.0120\\
1000 人的模擬 & 0.672 & 0.222 & 0.060 & 0.028 & 0.018\\
\bottomrule
\end{tabular}}

\newcommand{\ThreeFigFoot}{%
\begin{tikzpicture}[x=2.6cm,y=3.7cm,>=Stealth]
\draw[->] (1,0) -- (4.15,0) node[right,font=\scriptsize] {$n$};
\draw[->] (1,0) -- (1,1.05) node[above,font=\scriptsize] {$\delta_n/n$};
\draw[cBar,very thick] plot coordinates {(1.0000,0.5168) (1.0414,0.4922) (1.0792,0.4700) (1.1139,0.4512) (1.1461,0.4339) (1.1761,0.4189) (1.2041,0.4050) (1.2304,0.3928) (1.2553,0.3812) (1.3010,0.3612) (1.3424,0.3440) (1.3617,0.3364) (1.3979,0.3224) (1.4314,0.3100) (1.4624,0.2989) (1.5051,0.2843) (1.5315,0.2757) (1.5682,0.2641) (1.6021,0.2539) (1.6335,0.2448) (1.6628,0.2366) (1.6990,0.2268) (1.7324,0.2182) (1.7634,0.2104) (1.7993,0.2018) (1.8325,0.1942) (1.8692,0.1861) (1.8976,0.1801) (1.9345,0.1726) (1.9685,0.1659) (2.0000,0.1600) (2.0334,0.1539) (2.0682,0.1478) (2.1004,0.1424) (2.1335,0.1371) (2.1673,0.1318) (2.1987,0.1272) (2.2330,0.1222) (2.2672,0.1175) (2.3010,0.1130) (2.3324,0.1090) (2.3674,0.1047) (2.3997,0.1008) (2.4330,0.0970) (2.4669,0.0933) (2.4997,0.0898) (2.5328,0.0865) (2.5670,0.0831) (2.5999,0.0800) (2.6335,0.0770) (2.6665,0.0741) (2.6998,0.0713) (2.7332,0.0686) (2.7664,0.0661) (2.8000,0.0636) (2.8331,0.0612) (2.8669,0.0588) (2.8998,0.0566) (2.9335,0.0545) (2.9666,0.0525) (3.0000,0.0505) (3.0334,0.0486) (3.0667,0.0467) (3.1000,0.0450) (3.1332,0.0433) (3.1667,0.0417) (3.2000,0.0401) (3.2333,0.0386) (3.2667,0.0371) (3.2999,0.0357) (3.3332,0.0344) (3.3666,0.0331) (3.4000,0.0318) (3.4333,0.0306) (3.4667,0.0295) (3.5000,0.0284) (3.5334,0.0273) (3.5667,0.0263) (3.6000,0.0253) (3.6334,0.0243) (3.6667,0.0234) (3.7000,0.0225) (3.7334,0.0217) (3.7666,0.0209) (3.8000,0.0201) (3.8333,0.0193) (3.8666,0.0186) (3.9000,0.0179) (3.9333,0.0172) (3.9667,0.0166) (4.0000,0.0160)};
\draw[cGauss,very thick] plot coordinates {(1.0000,0.7881) (1.0414,0.7691) (1.0792,0.7520) (1.1139,0.7374) (1.1461,0.7240) (1.1761,0.7123) (1.2041,0.7015) (1.2304,0.6919) (1.2553,0.6828) (1.3010,0.6670) (1.3424,0.6533) (1.3617,0.6472) (1.3979,0.6360) (1.4314,0.6259) (1.4624,0.6169) (1.5051,0.6049) (1.5315,0.5978) (1.5682,0.5882) (1.6021,0.5795) (1.6335,0.5718) (1.6628,0.5647) (1.6990,0.5563) (1.7324,0.5487) (1.7634,0.5419) (1.7993,0.5342) (1.8325,0.5273) (1.8692,0.5198) (1.8976,0.5142) (1.9345,0.5071) (1.9685,0.5008) (2.0000,0.4950) (2.0334,0.4891) (2.0682,0.4831) (2.1004,0.4776) (2.1335,0.4721) (2.1673,0.4667) (2.1987,0.4617) (2.2330,0.4564) (2.2672,0.4513) (2.3010,0.4463) (2.3324,0.4418) (2.3674,0.4368) (2.3997,0.4324) (2.4330,0.4279) (2.4669,0.4234) (2.4997,0.4191) (2.5328,0.4150) (2.5670,0.4107) (2.5999,0.4067) (2.6335,0.4027) (2.6665,0.3988) (2.6998,0.3950) (2.7332,0.3912) (2.7664,0.3876) (2.8000,0.3839) (2.8331,0.3804) (2.8669,0.3769) (2.8998,0.3735) (2.9335,0.3701) (2.9666,0.3668) (3.0000,0.3636) (3.0334,0.3604) (3.0667,0.3573) (3.1000,0.3542) (3.1332,0.3512) (3.1667,0.3482) (3.2000,0.3453) (3.2333,0.3424) (3.2667,0.3396) (3.2999,0.3368) (3.3332,0.3341) (3.3666,0.3314) (3.4000,0.3288) (3.4333,0.3262) (3.4667,0.3236) (3.5000,0.3211) (3.5334,0.3186) (3.5667,0.3162) (3.6000,0.3137) (3.6334,0.3114) (3.6667,0.3090) (3.7000,0.3067) (3.7334,0.3045) (3.7666,0.3022) (3.8000,0.3000) (3.8333,0.2979) (3.8666,0.2957) (3.9000,0.2936) (3.9333,0.2915) (3.9667,0.2895) (4.0000,0.2875)};
\draw[cExact,very thick] plot coordinates {(1.0000,0.9051) (1.0414,0.8908) (1.0792,0.8782) (1.1139,0.8676) (1.1461,0.8581) (1.1761,0.8498) (1.2041,0.8423) (1.2304,0.8357) (1.2553,0.8296) (1.3010,0.8191) (1.3424,0.8102) (1.3617,0.8062) (1.3979,0.7991) (1.4314,0.7929) (1.4624,0.7874) (1.5051,0.7802) (1.5315,0.7760) (1.5682,0.7704) (1.6021,0.7655) (1.6335,0.7612) (1.6628,0.7574) (1.6990,0.7528) (1.7324,0.7488) (1.7634,0.7453) (1.7993,0.7414) (1.8325,0.7380) (1.8692,0.7344) (1.8976,0.7317) (1.9345,0.7284) (1.9685,0.7255) (2.0000,0.7229) (2.0334,0.7203) (2.0682,0.7177) (2.1004,0.7154) (2.1335,0.7131) (2.1673,0.7109) (2.1987,0.7089) (2.2330,0.7068) (2.2672,0.7049) (2.3010,0.7030) (2.3324,0.7014) (2.3674,0.6996) (2.3997,0.6981) (2.4330,0.6965) (2.4669,0.6950) (2.4997,0.6936) (2.5328,0.6923) (2.5670,0.6909) (2.5999,0.6897) (2.6335,0.6885) (2.6665,0.6874) (2.6998,0.6863) (2.7332,0.6852) (2.7664,0.6842) (2.8000,0.6832) (2.8331,0.6823) (2.8669,0.6814) (2.8998,0.6806) (2.9335,0.6797) (2.9666,0.6790) (3.0000,0.6782) (3.0334,0.6775) (3.0667,0.6768) (3.1000,0.6761) (3.1332,0.6755) (3.1667,0.6748) (3.2000,0.6742) (3.2333,0.6737) (3.2667,0.6731) (3.2999,0.6726) (3.3332,0.6721) (3.3666,0.6716) (3.4000,0.6711) (3.4333,0.6707) (3.4667,0.6703) (3.5000,0.6699) (3.5334,0.6695) (3.5667,0.6691) (3.6000,0.6687) (3.6334,0.6684) (3.6667,0.6680) (3.7000,0.6677) (3.7334,0.6674) (3.7666,0.6671) (3.8000,0.6668) (3.8333,0.6665) (3.8666,0.6662) (3.9000,0.6660) (3.9333,0.6657) (3.9667,0.6655) (4.0000,0.6653)};
\draw[cBar,dashed] plot coordinates {(1.0000,0.5046) (1.0414,0.4811) (1.0792,0.4607) (1.1139,0.4426) (1.1461,0.4265) (1.1761,0.4120) (1.2041,0.3989) (1.2304,0.3870) (1.2553,0.3761) (1.3010,0.3568) (1.3424,0.3402) (1.3617,0.3327) (1.3979,0.3192) (1.4314,0.3071) (1.4624,0.2963) (1.5051,0.2821) (1.5315,0.2737) (1.5682,0.2623) (1.6021,0.2523) (1.6335,0.2434) (1.6628,0.2353) (1.6990,0.2257) (1.7324,0.2172) (1.7634,0.2095) (1.7993,0.2010) (1.8325,0.1935) (1.8692,0.1855) (1.8976,0.1795) (1.9345,0.1721) (1.9685,0.1655) (2.0000,0.1596) (2.0334,0.1536) (2.0682,0.1475) (2.1004,0.1422) (2.1335,0.1368) (2.1673,0.1316) (2.1987,0.1270) (2.2330,0.1220) (2.2672,0.1173) (2.3010,0.1128) (2.3324,0.1088) (2.3674,0.1045) (2.3997,0.1007) (2.4330,0.0969) (2.4669,0.0932) (2.4997,0.0898) (2.5328,0.0864) (2.5670,0.0831) (2.5999,0.0800) (2.6335,0.0770) (2.6665,0.0741) (2.6998,0.0713) (2.7332,0.0686) (2.7664,0.0660) (2.8000,0.0635) (2.8331,0.0612) (2.8669,0.0588) (2.8998,0.0566) (2.9335,0.0545) (2.9666,0.0524) (3.0000,0.0505) (3.0334,0.0486) (3.0667,0.0467) (3.1000,0.0450) (3.1332,0.0433) (3.1667,0.0416) (3.2000,0.0401) (3.2333,0.0386) (3.2667,0.0371) (3.2999,0.0357) (3.3332,0.0344) (3.3666,0.0331) (3.4000,0.0318) (3.4333,0.0306) (3.4667,0.0295) (3.5000,0.0284) (3.5334,0.0273) (3.5667,0.0263) (3.6000,0.0253) (3.6334,0.0243) (3.6667,0.0234) (3.7000,0.0225) (3.7334,0.0217) (3.7666,0.0209) (3.8000,0.0201) (3.8333,0.0193) (3.8666,0.0186) (3.9000,0.0179) (3.9333,0.0172) (3.9667,0.0166) (4.0000,0.0160)};
\draw[cGauss,dashed] plot coordinates {(1.3979,0.9760) (1.4314,0.9532) (1.4624,0.9330) (1.5051,0.9065) (1.5315,0.8909) (1.5682,0.8700) (1.6021,0.8516) (1.6335,0.8353) (1.6628,0.8206) (1.6990,0.8031) (1.7324,0.7876) (1.7634,0.7737) (1.7993,0.7583) (1.8325,0.7445) (1.8692,0.7299) (1.8976,0.7190) (1.9345,0.7053) (1.9685,0.6931) (2.0000,0.6822) (2.0334,0.6710) (2.0682,0.6597) (2.1004,0.6496) (2.1335,0.6395) (2.1673,0.6295) (2.1987,0.6205) (2.2330,0.6110) (2.2672,0.6018) (2.3010,0.5929) (2.3324,0.5850) (2.3674,0.5763) (2.3997,0.5686) (2.4330,0.5608) (2.4669,0.5531) (2.4997,0.5458) (2.5328,0.5387) (2.5670,0.5315) (2.5999,0.5248) (2.6335,0.5181) (2.6665,0.5117) (2.6998,0.5054) (2.7332,0.4992) (2.7664,0.4932) (2.8000,0.4873) (2.8331,0.4816) (2.8669,0.4759) (2.8998,0.4705) (2.9335,0.4651) (2.9666,0.4599) (3.0000,0.4548) (3.0334,0.4498) (3.0667,0.4449) (3.1000,0.4401) (3.1332,0.4355) (3.1667,0.4308) (3.2000,0.4264) (3.2333,0.4220) (3.2667,0.4177) (3.2999,0.4135) (3.3332,0.4093) (3.3666,0.4053) (3.4000,0.4013) (3.4333,0.3974) (3.4667,0.3936) (3.5000,0.3898) (3.5334,0.3861) (3.5667,0.3825) (3.6000,0.3790) (3.6334,0.3755) (3.6667,0.3721) (3.7000,0.3687) (3.7334,0.3655) (3.7666,0.3622) (3.8000,0.3590) (3.8333,0.3559) (3.8666,0.3529) (3.9000,0.3498) (3.9333,0.3469) (3.9667,0.3440) (4.0000,0.3411)};
\draw[cGauss,dotted,thick] plot coordinates {(1.0000,0.7169) (1.0414,0.7017) (1.0792,0.6883) (1.1139,0.6764) (1.1461,0.6658) (1.1761,0.6562) (1.2041,0.6475) (1.2304,0.6395) (1.2553,0.6321) (1.3010,0.6190) (1.3424,0.6076) (1.3617,0.6024) (1.3979,0.5929) (1.4314,0.5845) (1.4624,0.5768) (1.5051,0.5665) (1.5315,0.5604) (1.5682,0.5521) (1.6021,0.5446) (1.6335,0.5379) (1.6628,0.5317) (1.6990,0.5243) (1.7324,0.5177) (1.7634,0.5117) (1.7993,0.5049) (1.8325,0.4987) (1.8692,0.4921) (1.8976,0.4872) (1.9345,0.4808) (1.9685,0.4751) (2.0000,0.4700) (2.0334,0.4646) (2.0682,0.4592) (2.1004,0.4543) (2.1335,0.4493) (2.1673,0.4444) (2.1987,0.4399) (2.2330,0.4351) (2.2672,0.4304) (2.3010,0.4258) (2.3324,0.4217) (2.3674,0.4172) (2.3997,0.4131) (2.4330,0.4090) (2.4669,0.4049) (2.4997,0.4010) (2.5328,0.3971) (2.5670,0.3932) (2.5999,0.3895) (2.6335,0.3858) (2.6665,0.3822) (2.6998,0.3787) (2.7332,0.3752) (2.7664,0.3718) (2.8000,0.3684) (2.8331,0.3652) (2.8669,0.3619) (2.8998,0.3588) (2.9335,0.3556) (2.9666,0.3526) (3.0000,0.3496) (3.0334,0.3466) (3.0667,0.3437) (3.1000,0.3408) (3.1332,0.3380) (3.1667,0.3352) (3.2000,0.3325) (3.2333,0.3299) (3.2667,0.3272) (3.2999,0.3246) (3.3332,0.3221) (3.3666,0.3195) (3.4000,0.3171) (3.4333,0.3146) (3.4667,0.3122) (3.5000,0.3099) (3.5334,0.3075) (3.5667,0.3052) (3.6000,0.3030) (3.6334,0.3008) (3.6667,0.2986) (3.7000,0.2964) (3.7334,0.2943) (3.7666,0.2922) (3.8000,0.2901) (3.8333,0.2881) (3.8666,0.2860) (3.9000,0.2841) (3.9333,0.2821) (3.9667,0.2802) (4.0000,0.2783)};
\draw[cExact,dashed] (1,0.6595) -- (4,0.6595);
\foreach \e in {1,2,3,4} {\draw (\e,-0.012) -- (\e,0.012); \node[below,font=\scriptsize] at (\e,0) {$10^{\e}$};}
\foreach \v in {0.2,0.4,0.6,0.8,1} {\draw (0.97,\v) -- (1.03,\v); \node[left,font=\scriptsize] at (0.97,\v) {\v};}
\node[cExact,font=\scriptsize,below] at (3.3,0.649) {三維 sc：$\to 1/G=0.6595$（虛線）};
\node[cGauss,font=\scriptsize,right] at (3.0,0.2) {二維：$\pi/\ln n$（虛線）、$\pi/\ln 8n$（點線）};
\node[cBar,font=\scriptsize,right] at (2.0,0.27) {一維：$\sqrt{8/\pi n}$（虛線）};
\end{tikzpicture}}

\newcommand{\ThreeTabFoot}{%
{\small\setlength{\tabcolsep}{4pt}\begin{tabular}{@{}rcccccccc@{}}
\toprule
 & \multicolumn{2}{c}{一維} & \multicolumn{3}{c}{二維（正方）} & \multicolumn{2}{c}{三維（sc）}\\
\cmidrule(lr){2-3}\cmidrule(lr){4-6}\cmidrule(lr){7-8}
$n$ & 精確 $\delta_n$ & $\sqrt{8n/\pi}$ & 精確 $\delta_n$ & $\pi n/\ln n$ & $\pi n/\ln 8n$ & 精確 $\delta_n$ & $0.6595n+0.574\sqrt n$\\
\midrule
$10^{1}$ & 5.17 & 5.05 & 7.9 & 13.6 & 7.2 & 9.1 & 8.4\\
$10^{2}$ & 16.00 & 15.96 & 49.5 & 68.2 & 47.0 & 72.3 & 71.7\\
$10^{3}$ & 50.48 & 50.46 & 363.6 & 454.8 & 349.6 & 678.2 & 677.6\\
$10^{4}$ & 159.58 & 159.58 & 2874.6 & 3410.9 & 2782.7 & 6652.6 & 6652.0\\
\bottomrule
\end{tabular}}}

\newcommand{\ThreeTabErrata}{%
\begin{tabular}{@{}lll@{}}
\toprule
位置 & 書上 & 應為\\
\midrule
p.~52--53，(3.3)(3.5) & {\color{cWrong}$\prod_{i=1}^{N}$} & {\color{cExact}$\prod_{i=0}^{N}$（起點也算腳印）}\\
p.~56，(3.18) 第二行 & {\color{cWrong}沒有 $\frac{1}{2\pi}$} & {\color{cExact}前面乘 $\frac{1}{2\pi}$}\\
p.~56，(3.19) 第二行 & {\color{cWrong}$\alpha^2-\alpha/2+1/z$} & {\color{cExact}$\alpha^2-\alpha/z+1$}\\
p.~57，(3.20) 後 & {\color{cWrong}$z<1$} & {\color{cExact}$0<z<\frac12$}\\
p.~57，(3.23) & {\color{cWrong}$z_c^{-N}N^{-1/2}/\Gamma(\frac12)$} & {\color{cExact}$z_c^{-N-1/2}N^{-1/2}/\Gamma(\frac12)$}\\
p.~57，(3.24) & {\color{cWrong}$N^{-1/2}/\Gamma(\frac12)=1/\sqrt{\pi N}$} & {\color{cExact}$\sqrt{2}N^{-1/2}/\Gamma(\frac12)=\sqrt{2/\pi N}$}\\
p.~59，(3.27) 後 & {\color{cWrong}sc 的 $c_2=1/2$} & {\color{cExact}$c_2=1$}\\
p.~60 & {\color{cWrong}$f''(z_c)=1/\delta^2$} & {\color{cExact}$f''(z_c-\delta)\approx1/\delta^2$}\\
p.~61，(3.32) 與前一句 & {\color{cWrong}$a$（$q=Qa(z_c-z)^{1/2}/c_2^{1/2}$）} & {\color{cExact}$c_1$}\\
p.~66，(3.53) 後 & {\color{cWrong}$(z_c-z)^{1/2}$ 的係數} & {\color{cExact}$(z_c-z)^{-1/2}$ 的係數}\\
p.~66，(3.56) & {\color{cWrong}$\sum_{n'=2}^{\infty}$} & {\color{cExact}$\sum_{n'=2}^{n}$}\\
p.~67，(3.57) & {\color{cWrong}積分下限 1} & {\color{cExact}積分下限 2}\\
\bottomrule
\end{tabular}}


\section{生成函數 II：回家、維度、走過幾個點（第 3 章）}

\begin{keybox}[先看這裡]
醉漢在街上亂走，遲早會走回家嗎？鳥在天上亂飛呢？這一章用第 2 章的生成函數回答兩個問題，答案取決於\textbf{維度}：
\begin{center}
\small
\begin{tabular}{@{}ll@{}}
回不回得了家 & 新的生成函數數「到訪幾次」；永遠不回家的機率 $1/G(z_c;0)$（3.2--3.3 節）\\
維度的角色 & 一維、二維遲早回家；三維不一定。分水嶺在 $d=2$（3.3--3.6 節）\\
走過幾個點 & 第 $n$ 步踩到新點的機率＝還沒回家的機率，足跡大小隨之而來（3.7 節）
\end{tabular}
\end{center}
\textbf{記號}：$Z$＝每一步可走的方向數，$z_c=1/Z$（一維 $\frac12$、正方格子 $\frac14$、簡單立方 sc $\frac16$）；$G(z;0)=\sum_NC(N;\vec0)z^N$；
$u_N=C(N;\vec0)/Z^N$＝第 $N$ 步剛好在原點的機率；$P_N$＝走了 $N$ 步都還沒回過家的機率；$\Delta_n$＝第 $n$ 步踩到新點的機率；
$\delta_n$＝$n$ 步平均踩過幾個不同的點。

\textbf{全章的主例}：第 1 章的十步醉漢。$N$ 步之內\textbf{從來沒回過家}的走法有
\[1,\ 2,\ 2,\ 4,\ 6,\ 12,\ 20,\ 40,\ 70,\ 140,\ 252\qquad(N=0,1,\dots,10)\]
條。最後一個 252 恰好等於 $\binom{10}{5}$＝第 10 步剛好在家的走法數（3.3 節）；3.7 節再用同一串數算「走過幾個點」。
\begin{center}
\small
\begin{tabular}{@{}lll@{\qquad}lll@{}}
\toprule
節 & 內容 & 動畫 & 節 & 內容 & 動畫\\
\midrule
3.1 & 醉漢與醉鳥 & W301 & 3.5 & 維度的分水嶺 & W305\\
3.2 & 數到訪次數的生成函數 & W302 & 3.6 & 三種立方格子 & W306\\
3.3 & 一維：一定回家 & W303 & 3.7 & 走過幾個點 & W307--W308\\
3.4 & 二維：慢得驚人 & W304 & 3.8--3.9 & 練習解答、筆誤 & \\
\bottomrule
\end{tabular}
\end{center}
數字一律由 \texttt{codes/utils/walk\_rec.py} 重算，圖表由 \texttt{ch3/make\_figs.py} 產生並先做數字檢查；逐式校對見 \texttt{ref/md/ch03\_formulas.md}。
連續播放：\texttt{Ch3Demo}（W301--W308）。
\end{keybox}

% =====================================================================
\subsection{醉漢與醉鳥：Pólya 的問題}
\bk{pp.~51, 55}\ \anim{W301}

數學家角谷靜夫有一句名言：「醉漢終究會走回家，醉鳥卻可能永遠迷路。」在一維的街道或二維的棋盤街道上亂走，遲早會回到出發點；
在三維空間裡亂飛，卻有一定的機率永遠回不來。遲早一定回來的漫步叫\textbf{常返}（recurrent），不一定回來的叫\textbf{暫態}（transient）。

\textbf{最小的例子}：一維走 4 步，共 $2^4=16$ 條路線。第 1--4 步之中碰過原點的有 10 條；從沒回家的有 6 條：
$++++$、$+++-$、$++-+$ 和它們的鏡像。

\textbf{大一點的實驗}：一、二、三維（直線、正方格子、簡單立方格子）各讓 1000 個人走 1000 步（固定亂數種子）。
「理論」一欄是 $1000\times(1-P_N)$，由 3.3--3.5 節的方法精確算出。
\begin{table}[H]
\centering
\small
\ThreeTabSim
\caption{各 1000 人的模擬：回過家的人數。一、二維的比例一直往 1 爬（二維爬得很慢），三維卡在三分之一左右}\label{three-tab-sim}
\end{table}

\begin{thmbox}[Pólya 定理]
格子上的簡單隨機漫步：在一維與二維，遲早回到出發點的機率是 1（常返）；在三維簡單立方格子，回到出發點的機率只有
\[1-\frac{1}{G(z_c;0)}=1-\frac{1}{1.516386}=0.3405\qquad\text{（暫態）}.\]
\end{thmbox}

% =====================================================================
\subsection{數「到訪幾次」的生成函數}
\bk{pp.~51--54}\ \anim{W302}

\subsubsection{權重 \texorpdfstring{$(1-t)^M$}{(1-t) 的 M 次方}}
令 $C(N,M;\vec x,\vec y,\vec w)$＝從 $\vec x$ 走 $N$ 步到 $\vec y$、途中\textbf{到訪 $\vec w$ 恰好 $M$ 次}的走法數。腳印算第 $0,1,\dots,N$ 步，共 $N+1$ 個，所以 $M\le N+1$。
新的生成函數把每條走法乘上 $(1-t)^M$：
\[C'(N,t;\vec x,\vec y,\vec w)=\sum_{M=0}^{\infty}C(N,M;\vec x,\vec y,\vec w)(1-t)^M,\qquad
C(N,M;\cdot)=(-1)^M\frac{1}{M!}\frac{d^M}{dt^M}C'(N,t;\cdot)\Big|_{t=1}\tag{3.1, 3.2}\]
注意展開點是 $t=1$，不是 $0$。$(1-t)^M$ 從哪裡來？給每個腳印一個權重 $w(s_i)=1-t\,\delta_{s_i,S}$（踩在特別的點 $S$ 上才不是 1），
整條走法的權重是乘積 $\prod_i(1-t\,\delta_{s_i,S})=(1-t)^M$（(3.3)--(3.5)）。$t=1$ 時，踩過 $S$ 的走法權重全變成 0。

\begin{notebox}[〔注〕p.~52--53，(3.3)(3.5)]
書上的乘積寫成 $\prod_{i=1}^{N}$。要和「$M\le N+1$」以及 (3.8) 允許 $N_1=0$ 一致，應從 $i=0$ 起：起點也算一個腳印。
\end{notebox}

\textbf{例}：一維、$N=4$、從 $0$ 出發、$\vec w=+1$、終點不限。第 4 步一定在偶數點，所以只看前 3 步、第 4 步任意（乘 2）：
到訪 0 次的前 3 步是 $-+-$、$--+$、$---$（$3\times2=6$ 條），1 次是 $+++$、$+--$、$-++$（6 條），2 次是 $++-$、$+-+$（$2\times2=4$ 條）。
所以 $C'(4,t)=6+6(1-t)+4(1-t)^2=16-14t+4t^2$。用 (3.2) 取回：$C'(4,1)=6$；$-\frac{d}{dt}C'\big|_{t=1}=-(-14+8)=6$；$\frac12\frac{d^2}{dt^2}C'\big|_{t=1}=4$。

\subsubsection{對 \texorpdfstring{$t$}{t} 展開：排容}
把權重對 $t$ 展開：一次項 $-t\sum_i\delta_{s_i,S}$（(3.6)）、二次項 $t^2\sum_{i>j}\delta_{s_i,S}\delta_{s_j,S}$（(3.7)）……。乘上走法再加總，
$(-t)^k$ 的係數＝「挑出 $k$ 次到訪、其餘隨便」的總數，也就是把走法在第 $k$ 次挑出的到訪處切開、各段自由的卷積：
\[C(N;\vec x,\vec y)-t\!\!\sum_{N_1+N_2=N}\!\!C(N_1;\vec x,\vec w)C(N_2;\vec w,\vec y)
+t^2\!\!\!\!\sum_{\substack{N_1+N_2+N_3=N\\N_2\ge1}}\!\!\!\!C(N_1;\vec x,\vec w)C(N_2;\vec w,\vec w)C(N_3;\vec w,\vec y)-\cdots\tag{3.8}\]
例子裡：$-t$ 的係數 $14$＝「第 $N_1$ 步到 $+1$、之後隨便走」：$N_1=1$ 有 $1\cdot2^3=8$ 條、$N_1=3$ 有 $3\cdot2=6$ 條（也就是每次到訪各算一次：$6\cdot1+4\cdot2=14$）；
$t^2$ 的係數 $4$＝「到 $+1$、走 2 步再回到 $+1$、再隨便走 1 步」$=1\cdot2\cdot2$。$N_2\ge1$ 是為了不把「走 0 步又到訪」算進去。
\begin{figure}[H]
\centering
\ThreeFigVirial
\caption{書上 Fig.~3.1 的樣子：(3.8) 的各項。紅點是 $\vec w$；$-t$ 項經過它一次，$+t^2$ 項在它那裡繞一圈再出發……}\label{three-fig-virial}
\end{figure}

\subsubsection{乘上 \texorpdfstring{$z^N$}{z 的 N 次方}：幾何級數}
乘 $z^N$ 對 $N$ 求和，卷積變成乘積；$t^k$ 項是 $G_{xw}G_1^{\,k-1}G_{wy}$（$G_{xw}=G(z;\vec x,\vec w)$ 等；$G_1$＝至少走一步的 $G(z;\vec w,\vec w)$），是公比 $-tG_1$ 的幾何級數：
\[G(z,t;\vec x,\vec y,\vec w)=G(z;\vec x,\vec y)-t\,\frac{G(z;\vec x,\vec w)\,G(z;\vec w,\vec y)}{1+t\,G_1(z;\vec w,\vec w)}.\tag{3.9}\]
(3.10) 說：對 $t$ 微分 $M$ 次、乘 $(-1)^M/M!$、令 $t=1$，就是「到訪恰好 $M$ 次」的生成函數。實際算出來（令 $s=1-t$ 展開）：
\[M=0:\ G_{xy}-\frac{G_{xw}G_{wy}}{G_{ww}},\qquad M\ge1:\ \frac{G_{xw}}{G_{ww}}\Big(\frac{G_{ww}-1}{G_{ww}}\Big)^{M-1}\frac{G_{wy}}{G_{ww}}.\]
讀法：$\frac{G_{xw}}{G_{ww}}$＝第一次到達 $\vec w$；$\frac{G_{ww}-1}{G_{ww}}$＝「離開再第一次回來」，出現 $M-1$ 次；最後離開 $\vec w$ 不再回來走到 $\vec y$。
（\texttt{walk\_rec} 用這個式子和暴力列舉互相核對到 $N=8$。）

% =====================================================================
\subsection{一維：一定回家}
\bk{pp.~55--58}\ \anim{W303}

\subsubsection{從不回家的生成函數}
令 $\vec w=\vec x$。為了不把第 0 步算成「回家」，分子用 $G_1$：
\[\frac{G(z;\vec x-\vec y)}{1+G_1(z;0)}=\frac{G(z;\vec x-\vec y)}{G(z;0)}\tag{3.12}\]
是「走到 $\vec y$、途中沒回過 $\vec x$」的生成函數。對終點 $\vec y$ 求和，用 $\sum_{\vec y}G(z;\vec x-\vec y)=g(z;\vec q=0)=\frac{1}{1-z\chi(0)}=\frac{z_c}{z_c-z}$（(3.14)(3.15)）：
\[\boxed{\ \text{$N$ 步都不回家的走法數}=[z^N]\ \frac{z_c}{z_c-z}\,\frac{1}{G(z;0)}\ }\tag{3.16}\]

\begin{notebox}[另一個看法（機率版）]
「第 $N$ 步在家」的路線＝「第 $k$ 步第一次回家」接上「從家出發、$N-k$ 步後在家」：$u_N=\sum_{k=1}^{N}f_ku_{N-k}$（$N\ge1$）。
令 $U(s)=\sum u_Ns^N$、$F(s)=\sum f_Ns^N$，得 $U=1+FU$，即 $1-F=1/U$。「$N$ 步都還沒回家」的機率 $P_N=1-\sum_{k\le N}f_k=[s^N]\frac{1-F(s)}{1-s}=[s^N]\frac{1}{(1-s)U(s)}$。
令 $s=z/z_c$，$U(s)=G(z;0)$，這就是 (3.16) 除以 $Z^N$。$N\to\infty$：$P_\infty=1-F(1)=1/U(1)=1/G(z_c;0)$。
\end{notebox}

\subsubsection{一維算到底}
一維 $\chi(q)=e^{iq}+e^{-iq}$，$z_c=\frac12$（(3.17)）。令 $\alpha=e^{iq}$，積分變成單位圓上的圍道積分：
\[G(z,0)=\frac{1}{2\pi}\int_0^{2\pi}\frac{dq}{1-z(e^{iq}+e^{-iq})}=-\frac{1}{2\pi iz}\oint\frac{d\alpha}{\alpha^2-\alpha/z+1}
=\frac{1}{\sqrt{1-4z^2}}=\frac{z_c}{\sqrt{z_c-z}\sqrt{z_c+z}}.\tag{3.18--3.21}\]
極點 $\alpha_\pm=\frac12\{\frac1z\pm\sqrt{(1/z)^2-4}\}$，$\alpha_+\alpha_-=1$，只有 $\alpha_-$ 在單位圓內（$0<z<\frac12$ 時）。
\begin{errbox}[p.~56--57，(3.18)--(3.20)]
(3.18) 第二行漏了 $\frac{1}{2\pi}$。(3.19) 第二行的分母印成 $\alpha^2-\alpha/2+1/z$，應為 $\alpha^2-\alpha/z+1$（它的兩根才是 (3.20) 的 $\alpha_\pm$）。
(3.20) 後說「$z<1$ 時只有 $\alpha_-$ 在圓內」：應為 $0<z<\frac12$；$\frac12<z<1$ 時根號裡是負的，兩個極點都在單位圓上。
\end{errbox}

$G(z;0)=\frac{1}{\sqrt{1-4z^2}}=\sum_k\binom{2k}{k}z^{2k}=1+2z^2+6z^4+20z^6+70z^8+252z^{10}+\cdots$，正是第 $N$ 步在家的走法數。
(3.16) 在一維可以精確算：
\begin{align*}
\frac{g(z;0)}{G(z;0)}&=\frac{\sqrt{1-4z^2}}{1-2z}=\sqrt{\frac{1+2z}{1-2z}}=(1+2z)\sum_k\binom{2k}{k}z^{2k}\\
&=1+2z+2z^2+4z^3+6z^4+12z^5+\cdots+252z^{10}+\cdots
\end{align*}
所以 $N=2k$ 步都不回家的走法數＝$\binom{2k}{k}$＝第 $2k$ 步剛好在家的走法數；奇數步 $2k+1$ 是它的 2 倍。
\begin{table}[H]
\centering
\ThreeTabNever
\caption{一維：$N$ 步都不回家的走法數（$\sqrt{(1+2z)/(1-2z)}$ 的係數）與機率 $P_N$。偶數 $N$ 那一欄和「第 $N$ 步在家」一樣多}\label{three-tab-never}
\end{table}

\subsubsection{\texorpdfstring{$N$ 很大：$P_N\approx\sqrt{2/\pi N}$}{N 很大：P ≈ √(2/πN)}}
$z\to z_c$ 時 $\sqrt{z_c+z}\to\sqrt{2z_c}=1$，(3.22) 給 $g/G\to(z_c-z)^{-1/2}$。補充 3 的 (2:S3-10)：$(z_c-z)^{-\alpha}$ 的 $z^N$ 係數約 $z_c^{-N-\alpha}N^{\alpha-1}/\Gamma(\alpha)$，
取 $\alpha=\frac12$ 再除以全部的 $z_c^{-N}=2^N$：
\[P_N\approx z_c^{-1/2}\frac{N^{-1/2}}{\Gamma(1/2)}=\sqrt{\frac{2}{\pi N}}\ \to\ 0.\tag{3.23, 3.24}\]
\begin{errbox}[p.~57，(3.23)(3.24)]
書上 (3.23) 寫 $\frac{z_c^{-N}}{\Gamma(1/2)}N^{-1/2}$、(3.24) 寫 $N^{-1/2}\frac{1}{\Gamma(1/2)}=\frac{1}{\sqrt{\pi N}}$：都漏了 $z_c^{-1/2}=\sqrt2\approx1.4142$。
精確值站在更正這邊（$N=10$、100、1000）：精確 $0.2461$、$0.0796$、$0.0252$；更正式 $0.2523$、$0.0798$、$0.0252$；
書上印的 $1/\sqrt{\pi N}$ 只有 $0.1784$、$0.0564$、$0.0178$。書上自己的 (3.54) 也用 $\sqrt{2/\pi}$。
\end{errbox}

\subsubsection{偏硬幣（練習 3.2--3.4）}
向右的機率 $p$、向左 $q=1-p$。用機率的生成函數（$z_c=1$）：$G(z;0)=\sum_k\binom{2k}{k}(pq)^kz^{2k}=1/\sqrt{1-4pqz^2}$，於是
\[P_\infty=\frac{1}{G(1;0)}=\sqrt{1-4pq}=\sqrt{(p+q)^2-4pq}=|p-q|.\]
$p=0.6$：$G(1;0)=1/\sqrt{0.04}=5$，永遠不回家的機率 $0.2$；有限 $N$（$=2,4,10,20,50,100$）時 $P_N=0.52$、$0.4048$、$0.2940$、$0.2427$、$0.2094$、$0.2016$，很快逼近 0.2。
回家的機率 $1-|p-q|$（$p=0.6$ 時 0.8）；$p=1$ 時永遠不回家的機率是 1。只要 $p\ne\frac12$，就有一定比例的人一去不回。

% =====================================================================
\subsection{二維：回家，但慢得驚人}
\bk{pp.~58--60}\ \anim{W304}

正方格子回到原點的走法 $C(N;\vec0)=1,4,36,400,4900,63504$（$N=0,2,\dots,10$，第 2 章的主例）。用上面的「機率版」：
第一次回家的機率 $f_2,f_4,\dots,f_{10}=0.25,\ 0.0781,\ 0.0430,\ 0.0286,\ 0.0210$（$f_2=4/16$：走出去馬上走回來；$f_4=20/256$），
和 $0.4208$，所以 $P_{10}=0.5792$：十步之後還有一半以上的人沒回過家。

一般維度：(3.25) 要的是 $\frac{z_c}{(z_c-z)G(z;0)}$ 的係數，而 $G(z;0)=\frac{1}{(2\pi)^d}\int\frac{d^dq}{1-z\chi(\vec q)}$。$z$ 靠近 $z_c$ 時，
積分主要來自 $\vec q$ 很小的地方，那裡 $\chi(\vec q)=c_1-c_2|\vec q|^2+\cdots$（(3.27)），$c_1=2d$：
\[G(z;0)\approx\frac{K_d}{(2\pi)^d}\int\frac{1}{c_1}\frac{q^{d-1}dq}{(z_c-z)+c_2q^2/c_1^2},\qquad K_d=\frac{2\pi^{d/2}}{\Gamma(d/2)}.\tag{3.28, 3.29}\]
\begin{errbox}[p.~59，(3.27) 後]
書上說簡單立方格子 $c_2=1/2$。$\chi=2\sum_i\cos q_i=2d-|\vec q|^2+\cdots$，所以 $c_2=1$（3.5 節的 $B_1=9/\pi$ 也要 $c_2=1$ 才對得上數值）。
\end{errbox}
$d=2$ 時角度積分後是 $\int_0^{q_0}\frac{q\,dq}{z_c-z+c_2q^2/c_1^2}=\frac{c_1^2}{2c_2}\ln\frac{z_c-z+c_2q_0^2/c_1^2}{z_c-z}$（(3.30)）：$G(z;0)\propto-\ln(z_c-z)$，在 $z_c$ \textbf{對數發散}。
正方格子還有精確式 $G(z;0)=\frac2\pi K(4z)$（第 2 章的橢圓積分），$z\to z_c$ 時
\begin{gather*}
G(z;0)\approx\frac1\pi\Big[\ln8-\ln\Big(1-\frac{z}{z_c}\Big)\Big]\\
\Rightarrow\quad P_N\approx\frac{\pi}{\ln N}\ \text{（書上的首項）}\quad\text{或}\quad\frac{\pi}{\ln N+\ln8+\gamma}\ \text{（含常數項）}.
\end{gather*}
\begin{table}[H]
\centering
\ThreeTabNoReturn
\caption{三種維度「$N$ 步都還沒回家」的機率 $P_N$：精確值（更新方程）與漸近式。二維的 $\pi/\ln N$ 在 $N=10$ 時是 1.364，比 1 還大；它要 $N$ 非常大才準}\label{three-tab-noreturn}
\end{table}
\textbf{慢得驚人}：要 $\pi/\ln N=0.01$，得 $N=e^{100\pi}\approx10^{136}$ 步——比宇宙裡的原子數（約 $10^{80}$）還多。
回到表 \ref{three-tab-sim}：1000 步之內還沒回家的有 $1000-676=324$ 人，理論 $0.3218$。

\begin{notebox}[最陡下降（書上 Fig.~3.2）與一個小筆誤]
書上用最陡下降估 $\frac{1}{2\pi i}\oint\frac{1}{(z_c-z)\ln(z_c-z)}\frac{dz}{z^{N+1}}$：把被積函數寫成 $h(z)e^{f(z)}$，$h=1/\ln(z_c-z)$、$f=-\ln(z_c-z)-(N+1)\ln z$。
鞍點在 $z=z_c(1-1/N)=z_c-\delta$，路徑平行虛軸穿過它，得到常數 $\times z_c^{-(N+1)}/\ln N$。
〔注〕p.~60 寫「$f''(z_c)=1/\delta^2$」：$f''$ 在 $z_c$ 發散，應是鞍點上的 $f''(z_c-\delta)\approx1/\delta^2$。
\end{notebox}

% =====================================================================
\subsection{維度的分水嶺}
\bk{pp.~60--62}\ \anim{W305}

為什麼偏偏是 $d=2$？在 $z=z_c$，(3.28) 的被積函數在 $q\to0$ 附近是 $\frac{q^{d-1}}{q^2}=q^{d-3}$：
\begin{figure}[H]
\centering
\ThreeFigTail
\caption{$z=z_c$ 時 (3.28) 的被積函數 $q^{d-1}/q^2$。從 $\varepsilon$ 積到 1 的面積（$\varepsilon=0.1$、$0.01$、$0.001$）：
$d=1$ 是 9、99、999，$d=2$ 是 2.303、4.605、6.908（$=\ln\frac1\varepsilon$），$d=3$ 是 0.9、0.99、0.999（綠色，停在 1）}\label{three-fig-tail}
\end{figure}
所以 $d\le2$ 時 $G(z_c;0)=\infty$、$P_\infty=0$（常返）；$d=3$ 時 $G(z_c;0)$ 有限、$P_\infty>0$（暫態）。角度積分給的常數
$K_1=2$、$K_2=2\pi\approx6.283$、$K_3=4\pi\approx12.566$、$K_4=2\pi^2\approx19.739$ 只是 $d$ 維單位球的表面積（練習 3.5）。

\textbf{奇異的形狀}：(3.31) 把 $q^{d-1}=q^{d-3}\cdot q^2$ 拆開，再令 $q=Q\,c_1(z_c-z)^{1/2}/c_2^{1/2}$，得
\[G(z;0)\propto\frac{c_1^2}{c_2}\frac{q_0^{d-2}}{d-2}-\Big(\frac{c_1^2}{c_2}\Big)^{d/2}(z_c-z)^{(d-2)/2}\int_0^\infty\frac{Q^{d-3}dQ}{1+Q^2}+R\,(z_c-z).\tag{3.32}\]
$2<d<4$ 時第一項是有限的常數、第二項是 $(z_c-z)$ 的非整數正次方：$G(z_c;0)$ 有限，但導數在 $z_c$ 發散。
$d=1$ 換成 $(z_c-z)^{-1/2}$（3.3 節），$d=2$ 換成對數（3.4 節）。
\begin{errbox}[p.~61，(3.32) 與前一句]
換變數印成 $q=Qa(z_c-z)^{1/2}/c_2^{1/2}$、係數 $(a^2/c_2)^{d/2}$：$a$ 應為 $c_1$（$a$ 在本書是格距，這裡用不到）。
\end{errbox}
\begin{table}[H]
\centering
\small
\ThreeTabGzero
\caption{$z$ 逼近 $z_c$ 時的 $G(z;0)$：一維、二維一路發散，三維停在 1.516（第三列用 Bessel 積分算，並和級數 $\sum u_Nx^N$ 核對）}\label{three-tab-gzero}
\end{table}

$d=3$ 時 $G(z;0)\approx G(z_c;0)$，(3.33) 的係數是 $z_c^{-N}/G(z_c;0)$，除以全部的 $z_c^{-N}$：
\[P_\infty=\frac{1}{G(z_c;0)}\qquad(\text{sc：}1/1.516386=0.659463).\tag{3.33, 3.34}\]
\textbf{有限 $N$ 的修正（練習 3.6）}：$G(z;0)\approx B_0-B_1(z_c-z)^{1/2}$。代進 (3.16)：
$P_N\approx\frac{1}{B_0}+\frac{B_1\sqrt{z_c}}{B_0^2}\frac{1}{\sqrt{\pi N}}$。由 (3.28)(3.32)（$c_1=6$、$c_2=1$、$\int_0^\infty\frac{dQ}{1+Q^2}=\frac\pi2$）得 $B_1=9/\pi\approx2.865$，所以
\[P_N\approx0.659463+\frac{0.28696}{\sqrt N}\]
（$N=100$：0.688159，精確 0.688062；$N=1000$：0.668537，精確 0.668534）。

% =====================================================================
\subsection{三種立方格子：Watson 積分}
\bk{pp.~62--63}\ \anim{W306}

三維的答案取決於 $G(z_c;0)=\frac{1}{(2\pi)^3}\int\frac{d^3q}{1-\chi(\vec q)/\chi(0)}$（(3.35)），稱為 \textbf{Watson 積分}（1939）。三種立方格子：
\begin{gather*}
\chi_{sc}=2(\cos q_x+\cos q_y+\cos q_z),\qquad \chi_{bcc}=8\cos q_x\cos q_y\cos q_z,\\
\chi_{fcc}=4(\cos q_x\cos q_y+\cos q_y\cos q_z+\cos q_x\cos q_z).\tag{3.36--3.38}
\end{gather*}
\begin{table}[H]
\centering
\ThreeTabLattices
\caption{三種立方格子的 Watson 積分（(3.39)--(3.41)）與逃脫機率。鄰居越多越容易逃走}\label{three-tab-lattices}
\end{table}
書上只給數值；閉式是 Gamma 函數。\texttt{walk\_rec} 另用兩種算法核對：sc 用 $\int_0^\infty e^{-t}I_0(t/3)^3dt$（$I_0$ 是 Bessel 函數；
來自 $\frac{1}{1-x}=\int_0^\infty e^{-t(1-x)}dt$ 與 $\frac{1}{2\pi}\int e^{a\cos q}dq=I_0(a)$），bcc 用 $\sum_k\big(\binom{2k}{k}/4^k\big)^3$（bcc 的回家次數是一維的三次方），
兩者都和閉式一致到小數第 10 位；fcc 用精確的回家次數級數加尾巴，差 $<10^{-3}$。（書上 $G_{bcc}$ 印 1.393\,203，是截斷值。）

\textbf{回家幾次（練習 3.8）}：sc 回家的機率 $R=1-0.659463=0.340537$。每次回到家，往後的漫步就像重新開始，所以剛好回家 $m$ 次的機率
\[f_m=R^m(1-R),\qquad \sum_mf_m=1,\qquad \sum_m m f_m=\frac{R}{1-R}=G-1=0.516386,\]
$m=0,\dots,4$ 依序是 $0.6595$、$0.2246$、$0.0765$、$0.0260$、$0.0089$：平均回家半次。
\begin{table}[H]
\centering
\small
\ThreeTabFm
\caption{sc：剛好回家 $m$ 次的機率。1000 步內的精確值由第一次回家的機率 $f_n$ 卷積算出；模擬是 1000 人各走 1000 步}\label{three-tab-fm}
\end{table}

% =====================================================================
\subsection{走過幾個不同的點}
\bk{pp.~63--68}\ \anim{W307--W308}

\subsubsection{新點的機率＝還沒回家的機率}
令 $\Delta_n$＝第 $n$ 步踩到\textbf{新點}（以前沒到過）的機率，$\Delta_0=1$。平均足跡就是把每一步的新點加起來：
\[\delta_n=\sum_{n'=0}^{n}\Delta_{n'},\qquad \Delta_n=\sum_{\vec y\ne\vec x}F_n(\vec x-\vec y)\tag{3.42, 3.43}\]
（$F_n(\vec x-\vec y)$＝第 $n$ 步\textbf{第一次}到 $\vec y$ 的機率）。\textbf{例}：一維 4 步的 16 條路線，數每一步是不是新點：
$\Delta_0,\dots,\Delta_4=1,\ 1,\ \frac12,\ \frac12,\ \frac38$，平均足跡 $\delta_0,\dots,\delta_4=1,\ 2,\ \frac52,\ 3,\ \frac{27}8$（16 條的足跡總和 $54=16\times\frac{27}{8}$）。

這串 $\Delta_n$ 和表 \ref{three-tab-never} 的 $P_n=1,1,\frac12,\frac12,\frac38$ 一模一樣。\textbf{為什麼？}把路線倒過來走：「第 $n$ 步踩到新點」＝「倒著走的人 $n$ 步都沒回到他的出發點」。
例如 $0\to1\to0\to1\to2$ 的第 4 步踩到新點 2；倒著走 $2\to1\to0\to1\to0$，從 2 出發沒再回到 2。每條路線倒過來仍是一條等機率的路線，所以 $\Delta_n=P_n$。

書上的算法（Fig.~3.3）：走到 B 的路線扣掉「之前到過 B」的，和 3.2 節一樣是幾何級數：
\[G(z;\vec r_B-\vec r_A)-G(z;\vec r_B-\vec r_A)G_1(z;0)+\cdots=\frac{G(z;\vec r_B-\vec r_A)}{1+G_1(z;0)},\qquad 1+G_1(z;0)=G(z;0).\tag{3.44--3.47}\]
對 B 求和得 $\frac{g(z;0)}{1+G_1(z;0)}=\frac{z_c}{z_c-z}\frac{1}{G(z;0)}$（(3.48)）——和 (3.16) 是\textbf{同一個函數}。
所以一維精確：$\delta_n=(2n+1)\binom{2k}{k}/4^k$（$n=2k$）、$2n\binom{2k}{k}/4^k$（$n=2k+1$）；二維暴力列舉 $4^n$ 條（$n\le5$）也和生成函數一致。

\subsubsection{足跡有多大}
\begin{itemize}
\item \textbf{一維}：$\Delta_n\approx\sqrt{2/\pi n}$，$\delta_n\approx\sqrt{\frac2\pi}\sum_{n'=1}^n(n')^{-1/2}\approx2\sqrt{\frac2\pi}\,n^{1/2}=\sqrt{8n/\pi}$（(3.53)(3.54)）。
      $n=100$：精確 15.997，$\sqrt{800/\pi}=15.958$。足跡只長成 $\sqrt n$，和第 1 章「離家 $\sqrt N$ 遠」一致。
\item \textbf{二維}：$\Delta_n\approx\pi/\ln n$，$\delta_n\approx\pi\sum_{n'=2}^{n}\frac{1}{\ln n'}\approx\frac{\pi n}{\ln n}$（(3.55)--(3.57)）。
      動畫 W308 的 400 步足跡踩過 159 個點，平均 $\delta_{400}=162.57$。
\item \textbf{三維}：(3.58) $\delta_n\approx n/G(z_c;0)$；sc 是 $0.6595\,n$：$\delta_n/n=0.9051,\ 0.7229,\ 0.6782$（$n=10,100,1000$）。
      加上次一項 $2\times0.28696\sqrt n=0.5739\sqrt n$（練習 3.11）：$n=1000$ 給 677.6，精確 678.2。
\end{itemize}
\begin{errbox}[p.~66--67，(3.53)--(3.57)]
(3.53) 後「$(z_c-z)^{1/2}$ 的 $z^n$ 係數」應為 $(z_c-z)^{-1/2}$。(3.56) 的和上限印成 $\infty$，應為 $n$（否則發散）。
(3.57) 的積分下限印成 1（三項都在 $n'=1$ 發散），照前一段的說明應為 2。
\end{errbox}
\begin{figure}[H]
\centering
\ThreeFigFoot
\caption{平均足跡占步數的比例 $\delta_n/n$（實線：精確）與漸近式（虛線、點線）。一維、二維的比例趨近 0，三維停在 $1/G(z_c;0)$}\label{three-fig-foot}
\end{figure}
\begin{table}[H]
\centering
\ThreeTabFoot
\caption{平均足跡 $\delta_n$：精確值與漸近式（二維的 $\pi n/\ln n$ 收斂得很慢）}\label{three-tab-foot}
\end{table}

% =====================================================================
\subsection{練習 3.1--3.12 解答}
{\small
\begin{enumerate}[label=\textbf{3.\arabic*},itemsep=0.05em]
\item 推 (3.9)。\quad 答：(3.8) 第 $k$ 項是 $k+1$ 段卷積（$\vec x\to\vec w$、$k-1$ 段 $\vec w\to\vec w$ 各至少一步、$\vec w\to\vec y$），乘 $z^N$ 求和得 $(-t)^kG_{xw}G_1^{k-1}G_{wy}$；對 $k\ge1$ 是公比 $-tG_1$ 的幾何級數，和為 $-\frac{tG_{xw}G_{wy}}{1+tG_1}$。
\item 偏硬幣不回家的機率不是 0。\quad 答：$P_\infty=1/G(1;0)=\sqrt{1-4pq}=|p-q|>0$（$p\ne\frac12$）。
\item $p=1$ 時逃脫機率是 1。\quad 答：$|p-q|=1$；由生成函數：$pq=0$ 使 $G(z;0)=1$，$P_\infty=1$。
\item 偏硬幣終究回家的機率。\quad 答：$1-|p-q|$；$p=0.6$ 時 0.8。
\item 推 (3.29)。\quad 答：$\int e^{-|\vec r|^2}d^dr=\big(\int e^{-x^2}dx\big)^d=\pi^{d/2}$；用球座標 $=K_d\int_0^\infty r^{d-1}e^{-r^2}dr=K_d\,\Gamma(d/2)/2$，
所以 $K_d=2\pi^{d/2}/\Gamma(d/2)$。
\item 三維有限 $N$ 的修正。\quad 答：見 3.5 節：$P_N\approx\frac1{B_0}+\frac{B_1\sqrt{z_c}}{B_0^2\sqrt{\pi N}}$，修正 $\propto N^{-1/2}$；sc 為 $0.659463+0.28696/\sqrt N$。
\item sc 逃脫機率的式子。\quad 答：$P_\infty=\Big[\frac{1}{(2\pi)^3}\int_{\rm BZ}\frac{d^3q}{1-\frac13(\cos q_x+\cos q_y+\cos q_z)}\Big]^{-1}=0.659463$。
\item 回家恰好 $m$ 次的比例 $f_m$。\quad 答：由 (3.10)，到訪 $m+1$ 次（含起點）的生成函數是 $g\,(G-1)^m/G^{m+1}$，$z\to z_c$ 時得 $f_m=R^m(1-R)$，$R=1-1/G(z_c;0)$；$\sum_mf_m=(1-R)\frac{1}{1-R}=1$。
\item bcc 的平均足跡。\quad 答：$\delta_N\approx N/G_{bcc}(z_c;0)=0.717770\,N$，$G_{bcc}=\frac{1}{(2\pi)^3}\int\frac{d^3q}{1-\cos q_x\cos q_y\cos q_z}$。
\item 偏硬幣的平均足跡。\quad 答：$\Delta_n\to|p-q|$，所以 $\delta_N\approx|p-q|N$（$p=0.6$：$0.2N$；$N=100$ 精確 24.95）：每步平均往前推進 $|p-q|$ 格，新點都在前緣。
\item sc 足跡的次一項。\quad 答：$\Delta_n\approx\frac1{B_0}+\frac{0.28696}{\sqrt n}$，加總得 $\delta_N\approx0.659463N+0.5739N^{1/2}$：次一項是 $N^{1/2}$。
\item sc 的平均足跡。\quad 答：$\delta_N\approx N/1.516386=0.659463\,N$。
\end{enumerate}}

% =====================================================================
\subsection{第 3 章的筆誤一覽}
給 2004 年版。每一條都用程式核對過（\texttt{walk\_rec.check()}、\texttt{check\_numbers\_ch3()}）；
逐式的 LaTeX 版見 \texttt{ref/md/ch03\_formulas.md}。其中 (3.3)(3.5) 的乘積與 $f''(z_c)$ 兩條是書上沒處理好的地方（文中用灰框〔注〕），其餘在文中用紅框。
\begin{center}
\footnotesize
\ThreeTabErrata
\end{center}
